%! AP Statistics — Unit 1, Lesson 1.3 — Guided Notes
\documentclass[10pt]{article}
\usepackage{apstats-article}
\usepackage{apstats-boxes}

%=============================================================================
\begin{document}

\pageheader{Unit 1, Lesson 1.3}{Guided Notes}
\namedateperiod

\vspace{0.12in}

% ── OBJECTIVE ────────────────────────────────────────────────────────────────
\begin{objectivebox}
\small By the end of this lesson, I will be able to\dots\\[4pt]
\writeline\\[1pt]
\writeline\\[1pt]
\writeline
\end{objectivebox}

\vspace{0.1in}

% ── VOCABULARY ───────────────────────────────────────────────────────────────
\begin{vocabbox}
\termblanklong{Frequency}
\termblanklong{Relative frequency}
\termblanklong{Frequency table}
\termblanklong{Relative frequency table}
\termblanklong{Distribution (of a categorical variable)}
\end{vocabbox}

\vspace{0.1in}

% ── HOOK ─────────────────────────────────────────────────────────────────────
\begin{hookbox}
\small\textbf{Which school has more walkers? --- Fill in the missing column.}

\vspace{8pt}
\renewcommand{\arraystretch}{1.6}
\begin{tabularx}{\linewidth}{@{} >{\bfseries}p{0.15\linewidth}
    C{0.12\linewidth} C{0.2\linewidth}
    C{0.12\linewidth} C{0.2\linewidth} @{}}
\rowcolor{sky}
\textbf{Mode} & \textbf{School A} & \textbf{Rel.\ Freq.\ A}
              & \textbf{School B} & \textbf{Rel.\ Freq.\ B} \\
\midrule
Bus  & 120 & \blank{2.2cm} & 30  & \blank{2.2cm} \\
Car  & 80  & \blank{2.2cm} & 50  & \blank{2.2cm} \\
Walk & 40  & \blank{2.2cm} & 120 & \blank{2.2cm} \\
\midrule
\rowcolor{sky}\textbf{Total} & 240 & \blank{2.2cm} & 200 & \blank{2.2cm} \\
\end{tabularx}

\vspace{10pt}
After computing relative frequencies, which school has more walkers \textit{proportionally}?\\[6pt]
\writeline\quad
What does this tell you about using counts alone for comparison?\\[6pt]
\writeline
\end{hookbox}

\vspace{0.1in}

% ── DIRECT INSTRUCTION ───────────────────────────────────────────────────────
\begin{notesbox}{Building a Frequency and Relative Frequency Table}
\small
\begin{multicols}{2}

\textbf{\textcolor{navy}{Step-by-Step Process}}\\[4pt]
\begin{enumerate}[leftmargin=1.4em, itemsep=3pt]
  \item List all \blank{2.5cm} in the first column.
  \item \blank{2.5cm} the data (use tally marks).
  \item Record the \blank{2cm} (count) in the second column.
  \item Divide each count by \blank{1.5cm} to get the relative frequency.
  \item \textbf{Check:} relative frequencies must sum to \blank{1cm}.
\end{enumerate}

\columnbreak

\textbf{\textcolor{navy}{Formula}}\\[4pt]
\[
  \text{Relative frequency} = \frac{\text{frequency}}{\blank{2cm}}
\]
\vspace{4pt}
Can be expressed as a \blank{2cm} (decimal)\\
or a \blank{2cm} (multiply by 100).\\[8pt]
\textbf{\textcolor{navy}{Quick error check:}}\\
Sum of all relative frequencies $= $ \blank{1cm}\\
Sum of all relative frequencies $= $ \blank{1.5cm}\%

\end{multicols}
\end{notesbox}

\vspace{0.1in}

% ── WORKED EXAMPLE ───────────────────────────────────────────────────────────
\begin{notesbox}{Worked Example --- Pet Ownership Survey ($n = 40$)}
\small

\noindent A sample of 40 households is surveyed about the type of pet they own most
(if any). The raw data are summarized below. Complete the table.

\vspace{8pt}
\renewcommand{\arraystretch}{1.8}
\begin{tabularx}{\linewidth}{@{} >{\bfseries}p{0.22\linewidth}
                                    C{0.16\linewidth}
                                    C{0.22\linewidth}
                                    C{0.22\linewidth} @{}}
\rowcolor{sky}
\textbf{Pet type} & \textbf{Tally} & \textbf{Frequency} & \textbf{Rel.\ Frequency} \\
\midrule
Dog    & \small||||||||||  & 12 & \blank{2.5cm} \\
Cat    & \small|||||||     & 8  & \blank{2.5cm} \\
Fish   & \small||||        & 6  & \blank{2.5cm} \\
Bird   & \small|||         & 4  & \blank{2.5cm} \\
None   & \small||||||||||  & 10 & \blank{2.5cm} \\
\midrule
\rowcolor{sky}\textbf{Total} & & \blank{2cm} & \blank{2.5cm} \\
\end{tabularx}

\vspace{10pt}
\textbf{Write one AP-style sentence interpreting the most common category:}\\[6pt]
\writeline\\[1pt]\writeline

\end{notesbox}

\newpage

% ── INTERPRETING IN CONTEXT ──────────────────────────────────────────────────
\begin{notesbox}{Describing a Distribution --- AP Language}
\small
\begin{multicols}{2}

\textbf{\textcolor{navy}{What to include:}}
\begin{itemize}[itemsep=3pt]
  \item The \blank{2cm} with the highest relative frequency (most common).
  \item The \blank{2cm} with the lowest relative frequency (least common).
  \item The \blank{2cm} of individuals in the sample (to give context).
  \item \blank{2.5cm} --- always include this.
\end{itemize}

\columnbreak

\textbf{\textcolor{navy}{Model sentence frame:}}\\[4pt]
\textit{``Of the \blank{1.5cm} [individuals] surveyed, the most common
[variable] was \blank{2cm} with \blank{1.5cm}\% of responses.
The least common was \blank{2cm} with \blank{1.5cm}\%.''}

\end{multicols}

\vspace{6pt}

\textbf{Practice:} Using the pet table from the worked example, write a complete
description of the distribution of pet type.\\[6pt]
\writeline\\[1pt]\writeline\\[1pt]\writeline

\end{notesbox}

\vspace{0.1in}

% ── GUIDED PRACTICE ──────────────────────────────────────────────────────────
\begin{practicebox}
\small
\textbf{Complete the relative frequency table. Then answer the questions.}

\vspace{6pt}
A survey of 25 AP Statistics students asks: ``How do you get to school?''
Results: Car~10, Bus~8, Walk~4, Bike~3.

\vspace{8pt}
\renewcommand{\arraystretch}{1.8}
\begin{tabularx}{\linewidth}{@{} >{\bfseries}p{0.22\linewidth}
                                    C{0.18\linewidth}
                                    C{0.22\linewidth} @{}}
\rowcolor{sky}
\textbf{Mode} & \textbf{Frequency} & \textbf{Rel.\ Frequency} \\
\midrule
Car  & 10 & \blank{2.5cm} \\
Bus  & 8  & \blank{2.5cm} \\
Walk & 4  & \blank{2.5cm} \\
Bike & 3  & \blank{2.5cm} \\
\midrule
\rowcolor{sky}\textbf{Total} & \blank{2cm} & \blank{2.5cm} \\
\end{tabularx}

\vspace{10pt}
\begin{enumerate}[leftmargin=1.5em, itemsep=6pt]
  \item What proportion of students walk to school? \blank{3cm}
  \item What percentage of students do \textit{not} take the bus?
        Show your work.\quad\blank{5cm}
  \item Write one complete sentence describing the most common mode of
        transportation in context.\\[6pt]
        \writeline\\[1pt]\writeline
\end{enumerate}

\end{practicebox}

\vspace{0.1in}

% ── SPIRAL / BIG IDEAS ────────────────────────────────────────────────────────
\begin{spiralbox}
\small
\begin{multicols}{2}

\textbf{This lesson connects forward to\dots}
\begin{itemize}[itemsep=2pt]
  \item \textbf{Lesson 1.4:} Bar charts and pie charts display the same distributions visually.
  \item \textbf{Unit 2:} Two-way tables extend this to two categorical variables simultaneously.
  \item \textbf{Unit 4:} Relative frequency $\approx$ probability when the sample is representative.
\end{itemize}

\columnbreak

\textbf{Big Idea --- Write in your own words:}\\
\textit{Why must we use relative frequency (not count) to fairly compare
two groups of different sizes?}\\[2pt]
\writeline\\[1pt]\writeline\\[1pt]\writeline

\end{multicols}
\end{spiralbox}

\vspace{0.1in}

\begin{notesbox}{Additional Notes}
\vspace{0pt}
\writeline\\\writeline\\\writeline\\\writeline\\\writeline\\\writeline
\end{notesbox}

\end{document}
